\documentclass[answers]{exam}
\usepackage{../../mypackages}
\usepackage{../../macros}

\SolutionEmphasis{\color{blue}}
\renewcommand{\solutiontitle}{\noindent}

\title{Interrogation N°2}
\author{N. Bancel}
\date{16 Octobre 2025}

\begin{document}

\textbf{Collège Lycée Suger}
\hfill
\textbf{Mathématiques} \\

\textbf{Année 2025-2026}
\hfill
\textbf{Classe de 6ème} \par

{\let\newpage\relax\maketitle}

\begin{center}
\textbf{\textcolor{red}{Durée : 1 heure. La calculatrice n'est pas autorisée}} \\
\textbf{\textcolor{red}{Une réponse donnée sans justification sera considérée comme fausse}} \\
Cette interrogation contient \numquestions\ questions, sur \numpages\ pages et est notée sur 20. \\
Une copie mal tenue sera pénalisée.

\end{center}

\section*{Notions et rappels du Chapitre 1}

\subsection*{Exercice 1 — Lecture et écriture des nombres (1.5 points)}

\begin{questions}
  \question[0.5] Écrire en lettres le nombre suivant : \quad $80\,431$.

  \begin{solution}
Le nombre $80\,431$ s'écrit en lettres : quatre-vingt mille quatre cent trente-et-un.
\end{solution}

  \question[1] Déterminer pour le nombre $2\,389,461$ :
  \begin{parts}
    \part Le nombre de dizaines.
    \part Le chiffre des centièmes.
    \part Le chiffre des centaines.
    \part La partie décimale
  \end{parts}
\end{questions}

\begin{solution}
\begin{parts}
    \part Le nombre de dizaines est $238$ (il y a deux façons de tomber sur ce résultat : $2\,389,461 \div 10 = 238,9461$ et la partie entière est $238$. L'autre méthode consiste à identifier la position du \textbf{chiffre des dizaines} : c'est 8, et de prendre le nombre qui est à gauche du 8, en incluant le 8).
    \part Le chiffre des centièmes est $6$ (dans $2\,389,461$, le chiffre des centièmes correspond au $6$ du nombre $2\,389,4\underline{6}1$).
    \part Le chiffre des centaines est $3$ (dans $2\,389,461$, le chiffre des centaines correspond au $3$ du nombre $2\,\underline{3}89,461$).
    \part La partie décimale est $0,461$ (la partie entière est $2\,389$).
\end{parts}
\end{solution}

\subsection*{Exercice 2 — Fractions décimales et écritures décimales (2 points)}

\begin{questions}
  \question[0.5] Ecrire sous la forme d'une seule et même fraction décimale $45,981$

\begin{solution}
  L'écriture sous forme de fraction décimale est une écriture avec \textbf{une seule fraction}
\[
45,981 = \frac{45981}{1000}
\]
\end{solution}

  \question[0.5] Ecrire sous la forme d'une décomposition par parties (entier + fraction décimale) $2,097$

\begin{solution}
Pour décomposer le nombre \(2,097\) en une partie entière et une partie fractionnaire :

1. Identifions la partie entière et la partie décimale :
   \[
   2,097 = 2 + 0,097
   \]

2. La partie \(2\) est l'entier.

3. La partie \(0,097\) peut être exprimée comme une fraction :
   \[
   0,097 = \frac{97}{1000}
   \]
  Ainsi, le nombre \(2,097\) peut être décomposé par parties en :
\[
2,097 = 2 + \frac{97}{1000}
\]
\end{solution}

  \question[0.5] Ecrire sous la forme d'une décomposition par rang ($... + \frac{...}{10} + \frac{...}{100} + \frac{...}{1000} + ...$) le nombre $34,0052$

\begin{solution}
Pour écrire le nombre \(34,0052\) sous la forme d'une décomposition par rang, nous devons identifier la contribution de chaque chiffre selon sa position.

1. La partie entière : \(34\)
    \[
    34 = 30 + 4 
    \]

2. La partie décimale est \(0,0052\), que nous décomposons par position des chiffres :
    \begin{addmargin}[4em]{1em}
    \begin{compactitem}
        \item 0 à la position des dixièmes : \(\frac{0}{10}\)
        \item 0 à la position des centièmes : \(\frac{0}{100}\)
        \item 5 à la position des millièmes : \(\frac{5}{1000}\)
        \item 2 à la position des dix-millièmes : \(\frac{2}{10000}\)
    \end{compactitem}
    \end{addmargin}

Le nombre \(34,0052\) se décompose en :
\[
34,0052 = 30 + 4 + \frac{0}{10} + \frac{0}{100} + \frac{5}{1000} + \frac{2}{10000}
\]

Ainsi, \(34,0052\) peut être écrit sous la forme d'une décomposition par rang comme suit :
\[
34,0052 = 30 + 4 + \frac{0}{10} + \frac{0}{100} + \frac{5}{1000} + \frac{2}{10000}
\]
Cette expression montre clairement la contribution de chaque chiffre en fonction de son rang. On peut évidemment simplifier : 
\[
34,0052 = 30 + 4 + \frac{5}{1000} + \frac{2}{10000}
\]
\end{solution}

  \question[0.5] Ecrire sous forme décimale $\frac{341}{1000}$


\begin{solution}
Pour écrire la fraction \(\frac{341}{1000}\) sous forme décimale, il suffit de comprendre que le dénominateur \(1000\) indique que nous avons affaire à un nombre de millièmes. Ainsi, la fraction s'exprime directement avec trois chiffres après la virgule.

\[
\frac{341}{1000} = 0.341
\]

Une autre méthode consiste à effectuer le calcul mental de la division \(341 \div 1000\), ce qui donne également \(0.341\).

\end{solution}

\end{questions}

\subsection*{Exercice 3 — Lecture sur une droite graduée (2 points)}

\begin{questions}
  \question[1] Donner l'abscisse de $T$ en écriture décimale. Justifier. La justification compte autant que la réponse.
\end{questions}

  \begin{center}
    \includegraphics[width=0.8\linewidth]{interro_2_01.jpg}
  \end{center}

  \begin{solution}
Pour déterminer l'abscisse de $T$ en écriture décimale, observons la graduation sur l'axe où se trouve $T$.
\begin{compactitem}
  \item On identifie les abscisses qui nous sont données : $2.9$ et $3.0$.
  \item On compte le nombre de "parts" qui séparent $2.9$ et $3.0$. Il y a 10 petites graduations entre ces deux grandes graduations.
  \item L'écart entre $2.9$ et $3.0$ est de $0.1$. Cet écart est divisé en 10 parties égales. Donc chaque petite graduation représente un incrément de $0.1 \div 10 = 0.01$.
  \item En comptant les petites graduations à partir de $2.9$ jusqu'à $T$, on trouve que $T$ est situé à 7 petites graduations après $2.9$.
  \item Mathématiquement, cela se traduit par :
  \[ T = 2.9 + 7 \times 0.01 = 2.9 + 0.07 = 2.97\]
  \item L'abscisse de $T$ est donc $T(2.97)$
\end{compactitem}
\end{solution}

\begin{questions}
  \question[1] Donner l'abscisse du point $W$ en écriture décimale. Justifier.
\end{questions}

    \begin{center}
    \includegraphics[width=0.8\linewidth]{interro_2_02.jpg}
  \end{center}

\begin{solution}
Pour déterminer l'abscisse de $W$ en écriture décimale, observons la graduation près de $6.41$.
\begin{compactitem}
  \item Les abscisses repères sont $6.41$ et $6.42$.
  \item Il y a 10 petites graduations entre $6.41$ et $6.42$. L'écart entre ces deux valeurs est $0.01$, donc chaque petite graduation vaut $0.01 \div 10 = 0.001$.
  \item $W$ est situé à 4 petites graduations après $6.41$.
  \item Donc, en valeur décimale :
  \[
  W = 6.41 + 4 \times 0.001 = 6.41 + 0.004 = 6.414.
  \]
\end{compactitem}
Ainsi, l'abscisse de $W$ est $6.414$.
\end{solution}


\subsection*{Exercice 4 — Comparaisons et rangement (2 points)}

\begin{questions}
  \question[0.5] Effectuer une décomposition du nombre $903\,821$
  \begin{solution}
Pour décomposer le nombre \(903\,821\), écrivons-le comme une somme où chaque chiffre est multiplié par la valeur de sa position.

\[
903\,821 = 9 \times 100\,000 + 0 \times 10\,000 + 3 \times 1\,000 + 8 \times 100 + 2 \times 10 + 1 \times 1
\]
\end{solution}
  \question[1.5] Classer les nombres suivants dans l'ordre \textbf{croissant} :  
  $87,09 \quad 82,2 \quad 82,92 \quad 87,764 \quad 82,467 \quad 82,76$.
\end{questions}

\begin{solution}
On aligne d'abord tous les nombres avec trois décimales pour pouvoir comparer rang par rang (des dizaines jusqu’aux millièmes) :
\[
\begin{array}{r}
82,200\\
82,467\\
82,76\\
82,920\\
87,090\\
87,764
\end{array}
\]

\begin{enumerate}
  \item \textbf{Rang des dizaines} : tous les nombres commencent par \(8\) → égalité, on passe au rang suivant.
  \item \textbf{Rang des unités} : on a \(2\) pour les \(82,xxxx\) et \(7\) pour les \(87,xxxx\).  
  Comme \(2 < 7\), tous les \(82,xxxx\) sont plus petits que tous les \(87,xxxx\).
  \item \textbf{À l’intérieur des \(82,xxxx\)} (on compare les chiffres après la virgule) :
    \begin{itemize}
      \item Rang des dixièmes : \(2 < 4 < 7 < 9\), donc  
      \[
      82,200 < 82,467 < 82,76 < 82,920
      \]
      (pas besoin d’aller plus loin, les dixièmes suffisent à départager).
    \end{itemize}
  \item \textbf{À l’intérieur des \(87,xxxx\)} :
    \begin{itemize}
      \item Rang des dixièmes : \(0 < 7\), donc  
      \[
      87,090 < 87,764
      \]
    \end{itemize}
\end{enumerate}

Ainsi, l’ordre \textbf{croissant} est :
\[
82,2 < 82,467 < 82,76 < 82,92 < 87,09 < 87,764
\]
\end{solution}

\section*{Applications du Chapitre 2}

\subsection*{Exercice 5 — Règles de priorité et calcul astucieux (5 points)}

\begin{questions}
  \question[2.5] Effectuer les calculs suivants en détaillant les étapes.
  \begin{parts}
    \part $A = 45 - (16 + 8) + 4$
    \part $B = 5 \times ((4 + 5) \times 2)$
    \part $C = (1 + 9) \times (14 - 3)$
    \part $D = 3 \times [10 - (2 + 3)] + 5$
  \end{parts}
\end{questions}

\begin{solution}
\begin{parts}

\part
\begin{align*}
A &= 45 - (16 + 8) + 4\\
  &= 45 - 24 + 4 && \text{on calcule la parenthèse}\\
  &= 21 + 4       && \text{car } 45 - 24 = 21\\
  &= 25
\end{align*}

\part
\begin{align*}
B &= 5 \times \big((4 + 5) \times 2\big)\\
  &= 5 \times (9 \times 2) && \text{on calcule } 4+5\\
  &= 5 \times 18\\
  &= 90
\end{align*}

\part
\begin{align*}
C &= (1 + 9) \times (14 - 3)\\
  &= 10 \times (14 - 3) && \text{on calcule } 1+9\\
  &= 10 \times 11       && \text{on calcule } 14-3\\
  &= 110
\end{align*}

\part
\begin{align*}
D &= 3 \times \big[10 - (2 + 3)\big] + 5\\
  &= 3 \times [10 - 5] + 5 && \text{on calcule } 2+3\\
  &= 3 \times 5 + 5\\
  &= 15 + 5\\
  &= 20
\end{align*}
\end{parts}

\end{solution}

\begin{questions}
  \question[2.5] Effectuer les calculs suivants en faisant des regroupements astucieux - Mettre en lumière ces groupements astucieux (je voudrais voir comment vous avez regroupé les nombres) :
  \begin{parts}
    \part $D = 25 + 18 + 75 + 82$
    \part $E = 4,5 \times 2 \times 5$ 
    \part $F = 2,5 \times 5 \times 4 \times 177 \times 2$
    \part $C = 25 \times 4,7 \times 0,5 \times 4 \times 2$
  \end{parts}
\end{questions}

\begin{solution}

  Avant de commencer les calculs, rappelons deux propriétés très importantes :  
la \textbf{commutativité} et l’\textbf{associativité}, qui s’appliquent à la fois à l’\textbf{addition} et à la \textbf{multiplication}.

\bigskip
\textbf{1. La commutativité}

\vspace{1em}
Cette propriété indique que l’ordre des nombres n’a pas d’importance :  
on peut changer leur place sans que le résultat change.

\[
a + b = b + a \qquad \text{et} \qquad a \times b = b \times a
\]

Exemples :
\[
3 + 5 = 5 + 3 = 8 \qquad \text{et} \qquad 4 \times 7 = 7 \times 4 = 28
\]

\bigskip
\textbf{2. L’associativité}
\vspace{1em}
Cette propriété indique que, lorsque plusieurs additions ou multiplications se suivent,  
on peut regrouper les nombres comme on veut sans changer le résultat.

\[
(a + b) + c = a + (b + c) \qquad \text{et} \qquad (a \times b) \times c = a \times (b \times c)
\]

Exemples :
\[
(2 + 3) + 5 = 2 + (3 + 5) = 10 \qquad \text{et} \qquad (2 \times 3) \times 5 = 2 \times (3 \times 5) = 30
\]

\bigskip
Ces deux propriétés nous permettent de \textbf{regrouper ou d’échanger les nombres} pour calculer plus facilement.

\bigskip


  \begin{parts}

\part
\begin{align*}
D &= 25 + 18 + 75 + 82\\
  &= (25 + 75) + (18 + 82) && \text{on regroupe pour faire des centaines}\\
  &= 100 + 100             && \text{car }25+75=100\ \text{et}\ 18+82=100\\
  &= 200
\end{align*}

\part
\begin{align*}
E &= 4,5 \times 2 \times 5\\
  &= 4,5 \times (2 \times 5) && \text{par associativité de la multiplication, on regroupe }2\ \text{et}\ 5\\
  &= 4,5 \times 10            && \text{car }2\times5=10\\
  &= 45                       && \text{multiplier par }10\ \text{déplace la virgule d’un rang}
\end{align*}

\part
\begin{align*}
F &= 2,5 \times 5 \times 4 \times 177 \times 2\\
  &= (2,5 \times 4) \times (5 \times 2) \times 177 && \text{commutativité/associativité pour créer des }10\\
  &= 10 \times 10 \times 177                       && \text{car }2,5\times4=10\ \text{et}\ 5\times2=10\\
  &= 100 \times 177\\
  &= 17\,700
\end{align*}

\part
\begin{align*}
C &= 25 \times 4,7 \times 0,5 \times 4 \times 2\\
  &= (25 \times 4) \times (0,5 \times 2) \times 4,7 && \text{on regroupe pour obtenir }100\ \text{et }1\\
  &= 100 \times 1 \times 4,7                       && \text{car }25\times4=100\ \text{et}\ 0,5\times2=1\\
  &= 100 \times 4,7\\
  &= 470
\end{align*}

\end{parts}
\end{solution}



\subsection*{Exercice 6 - Calculs posés. (3.5 points)}

\textbf{En posant le calcul, déterminer la valeur de :}
\begin{questions}
  \question[1] $G = 7496,8 - 29,54$
  \question[1.5] $H = 49,3 \times 2,05$
  \question[1] $I = 496,891 + 25,72$
\end{questions}

\begin{solution}
\begin{parts}

\part 
\[
G = 7496,8 - 29,54 = 7467,26
\]

\begin{center}
\includegraphics[width=0.5\linewidth]{interro_2/soustraction.jpg}
\end{center}

\part 
\[
H = 49,3 \times 2,05 = 101,065
\]

\begin{center}
\includegraphics[width=0.5\linewidth]{interro_2/multiplication.jpg}
\end{center}

\part 
\[
I = 496,891 + 25,72 = 522,611
\]


\begin{center}
\includegraphics[width=0.5\linewidth]{interro_2/addition.jpg}
\end{center}

\end{parts}
\end{solution}

\subsection*{Exercice 7 — Ordres de grandeur (1.5 point)}

\begin{questions}
  \question Donner un ordre de grandeur (et préciser à quel ordre de grandeur vous arrondissez) de :
  \begin{parts}
    \part[0.5] $J = 78,4 + 121,6$
    \part[0.5] $K = 2,98 \times 5,1$
    \part[0.5] $L = 1204,32 \times 98,8$
  \end{parts}

  \begin{solution}
\begin{parts}

\part 
\[
J = 78,4 + 121,6 \approx 80 + 120 = 200
\]
On arrondit à la dizaine la plus proche pour chaque nombre.  
L’ordre de grandeur de \(J\) est donc 200.

\part 
\[
K = 2,98 \times 5,1 \approx 3 \times 5 = 15
\]
On arrondit chaque facteur à l’unité la plus proche.  
L’ordre de grandeur de \(K\) est donc 15.

\part 
\[
L = 1204,32 \times 98,8 \approx 1200 \times 100 = 120000
\]
On arrondit à la centaine pour \(1204,32\) et à la centaine pour \(98,8\).  
L’ordre de grandeur de \(L\) est donc $120\,000$.

\end{parts}
\end{solution}
  
\end{questions}

\subsection*{Exercice 9 — Problème (2.5 point)}

\begin{questions}
  \question Je paye 100 feuilles à $0,05$ euros chacune, 1 classeur à 6 euros, et une règle à 2.4 euros. Je donne un billet de 50 euros.  
  Combien doit-on me rendre ?
\end{questions}

\begin{solution}

On cherche à déterminer la monnaie à rendre après l’achat.  
Commençons par calculer le prix total de tous les articles.

\bigskip
\textbf{1. Formule générale}
\[
Prix_{total} = Prix_{feuilles} + Prix_{classeur} + Prix_{regle}
\]

\bigskip
\textbf{2. Détail des calculs pour chaque article}

\begin{align*}
Prix_{feuilles} &= Prix_{une\_feuille} \times Nombre_{feuilles} && \text{formule du prix des feuilles}\\
                &= 0,05 \times 100 && \text{substitution des valeurs}\\
                &= 5\\[6pt]
Prix_{classeur} &= 6\\[6pt]
Prix_{regle} &= 2,4
\end{align*}

\bigskip
\textbf{3. Calcul du prix total}

\[
Prix_{total} = Prix_{feuilles} + Prix_{classeur} + Prix_{regle}
\]
\[
Prix_{total} = 5 + 6 + 2,4 = 13,4
\]

\bigskip
\textbf{4. Calcul de la monnaie à rendre}

\[
Monnaie_{rendue} = Billet_{donne} - Prix_{total}
\]
\[
Monnaie_{rendue} = 50 - 13,4 = 36,6
\]

\bigskip
\textbf{Conclusion :}  
On doit rendre au client une somme de \(36,6\)€.
\end{solution}

\end{document}
